% Section 1: Introduction
\subsection{Introduction}
\label{sec:forecasting_rl_intro}

Forecasting and reinforcement learning have grown up on adjacent benches. Time series
forecasters minimise mean-squared error or proper scoring rules. Reinforcement
learners maximise expected discounted reward. The two fields share an underlying
quantity, namely the conditional distribution of a future state, yet they disagree on
which functional of that distribution counts as evaluation. \citet{GrangerPesaran2000}
pointed out the consequence two decades ago. The same forecast can be more accurate
in a statistical sense and less useful in an economic sense whenever the decision
loss is asymmetric or threshold-driven.

This chapter surveys what happens at the intersection. The literature splits cleanly
into two halves and one bridge. In the first half, reinforcement learning is brought
\emph{to} forecasting, either as a meta-algorithm that selects ensemble weights, as a
mechanism that trains a predictor on a decision-aware reward, or as a tool that learns
value functions for sequential estimation. In the second half, forecasting is brought
\emph{to} reinforcement learning, in the form of world models, trajectory generative
models, counterfactual outcome predictors for off-policy evaluation, and the recent
generation of pretrained time series foundation models. The bridge concerns the
question that \citet{GrangerPesaran2000} and \citet{GneitingRaftery2007} formalised
from the statistical side, namely how the choice of loss governs what is identified.
Modern decision-focused learning, traceable to \citet{ElmachtoubGrigas2022} and
\citet{DontiAmosKolter2017}, makes the Wald (1950) statistical decision theory
differentiable, and the resulting machinery is now mature enough to deploy in
sequential settings.

We are economists, so we frame each piece by what it adds beyond ordinary least
squares with a plug-in decision rule. The chapter ends with one simulation that puts
the abstractions on a familiar substrate, namely an AR(1) demand process and a
single-period newsvendor decision, and a short list of open problems that an
economist is well placed to attack.
